% Part: first-order-logic
% Chapter: tableaux
% Section: derivations

\documentclass[../../../include/open-logic-section]{subfiles}

\begin{document}

\iftag{FOL}
      {\olfileid{fol}{tab}{der}}
      {\olfileid{pl}{tab}{der}}

\olsection{\usetoken{P}{tableau}}

\begin{explain}
اسیں دس چکے آں کہ مفروضہ کیہ اے، تے استنباط دے قاعدے وی دے چکے
آں۔ !!^{tableau}s ایہناں توں استقرائی طور تے بندے نیں: ہر
!!{tableau} یا تاں اک ایہو جیہی شاخ اے جیہڑی اک یا ودھ مفروضیاں
اُتے مشتمل ہووے، یا پھر شاخ اُتے استنباطی قاعدہ لا کے پہلے
!!a{tableau} توں حاصل ہوندا اے۔
\end{explain}

\begin{defn}[!!^{tableau}]
مفروضیاں $\sFmla{S_1}{!A_1}$، \dots،
$\sFmla{S_n}{!A_n}$ لئی !!^a{tableau} (جتھے ہر $S_i$ یا تاں
$\True$ اے یا~$\False$) !!{signed formula}s دا اک محدود درخت
اے جیہڑا ہِٹھ لکھیاں شرطاں پوری کردا اے:
\begin{enumerate}
\item درخت دے سب توں اُتّے والے $n$ عدد !!{signed formula}s،
  $\sFmla{S_i}{!A_i}$ نیں، اک دوجے دے ہِٹھاں۔
\item درخت وچ ہر اوہ !!{signed formula} جیہڑا مفروضیاں وچوں نہیں،
  اوہدے اُتّے والی شاخ دے کسے !!a{signed formula} اُتے استنباطی
  قاعدے دی درست ورتوں توں حاصل ہوندا اے۔
\end{enumerate}
!!a{tableau} دی شاخ اُدوں تے صرف اُدوں \emph{بند} اے جدوں اوہدے
وچ $\sFmla{\True}{!A}$ تے~$\sFmla{\False}{!A}$ دونیں
ہون، تے نہیں تاں \emph{کھلی} اے۔ جیہڑے !!^a{tableau} دی ہر شاخ
بند ہووے، اوہ (اپنے مفروضیاں دے سیٹ لئی) \emph{بند !!{tableau}}
اے۔ جے !!a{tableau} بند نہیں، یعنی اوہدے وچ گھٹ توں گھٹ اک کھلی
شاخ اے، تاں اوہ \emph{کھلا} اے۔
\end{defn}

\begin{ex}
  مفروضیاں دا ہر سیٹ اپنے آپ وچ !!a{tableau} اے، پر عام طور تے
  بند نہیں ہووے گا۔ (ظاہر اے کہ ایہ صرف اُدوں بند اے جدوں مفروضیاں
  وچ پہلے ای !!{signed formula}s دی جوڑی
  $\sFmla{\True}{!A}$ تے~$\sFmla{\False}{!A}$ موجود ہووے۔)

  % PNBTAB-003: bare A is the formula component of the signed formula; the sign remains present in the tableau.
  !!a{tableau} (کھلا یا بند) وچ موجود !!a{signed formula}، جیہدا
  فارمولا~$!A$ اے، اُتے استنباط دے قاعدیاں وچوں اک لا کے اس توں
  نواں، وڈا تابلو بنا سکدے آں۔ ایہ قاعدہ اوہدے لاگو ہون والی
  $!A$ دی آمد رکھن والی ہر شاخ دے سرے اُتے اک یا ودھ
  !!{signed formula}s شامل کرے گا۔

  مثال لئی، مفروضہ $\sFmla{\True}{!A \land \lnot
    !A}$ لو۔ صرف اوس مفروضے اُتے مشتمل (کھلا) !!{tableau} ایہ اے:
  \begin{center}
    \begin{tableau}{}
      [\sFmla{\True}{\formula{A} \land \lnot \formula{A}}, just=\TAss]
    \end{tableau}{}
  \end{center}
  ایس مفروضے اُتے $\TRule{\True}{\land}$ دا قاعدہ لا کے نواں
  !!{tableau} ملدا اے۔ ایہ قاعدہ سانوں !!{tableau} وچ دو نویاں
  لیکاں شامل کرن دی اجازت دیندا اے، $\sFmla{\True}{!A}$ تے
  $\sFmla{\True}{\lnot !A}$:
  \begin{center}
    \begin{tableau}{}
      [\sFmla{\True}{\formula{A} \land \lnot \formula{A}}, just=\TAss
        [\sFmla{\True}{\formula{A}}, just={\TRule{\True}{\land}[1]},
          [\sFmla{\True}{\lnot \formula{A}}, just={\TRule{\True}{\land}[1]}
          ]
        ]
      ]
    \end{tableau}{}
  \end{center}
  جدوں اسیں !!{tableau}s لکھدے آں تاں لاگو کیتے قاعدے سجے پاسے
  درج کردے آں (مثلاً $\TRule{\True}{\land} 1$ دا مطلب اے کہ
  اوس لکیر دا !!{signed formula}، $\TRule{\True}{\land}$ قاعدہ
  لکیر~$1$ دے !!{signed formula} اُتے لا کے
  حاصل ہویا اے)۔ ایس نویں !!{tableau} وچ ہور
  !!{signed formula}s نیں، پر صرف اک
  ($\sFmla{\True}{\lnot !A}$) اُتے اسیں قاعدہ لا سکدے آں
  (ایس صورت وچ $\TRule{\True}{\lnot}$ دا قاعدہ)۔ ایس توں بند
  !!{tableau} ملدا اے:
  \begin{center}
    \begin{tableau}{}
      [\sFmla{\True}{\formula{A} \land \lnot \formula{A}}, just=\TAss
        [\sFmla{\True}{\formula{A}}, just={\TRule{\True}{\land}[1]}
          [\sFmla{\True}{\lnot \formula{A}}, just={\TRule{\True}{\land}[1]}
            [\sFmla{\False}{\formula{A}}, just={\TRule{\True}{\lnot}[3]}, close]
          ]
        ]
      ]
    \end{tableau}{}
  \end{center}
\end{ex}

\end{document}
